\documentclass[a4paper,12pt]{article}

% ==============================
%            ПАКЕТЫ
% ==============================
\usepackage{cmap}
\usepackage[T2A]{fontenc}
\usepackage[utf8]{inputenc}
\usepackage[english,russian]{babel}
\usepackage{hyperref}
\usepackage{multicol}
\usepackage{mathtools}
\usepackage{amssymb}
\usepackage{enumitem}
\usepackage{tikz}
\usepackage{cancel}
\usetikzlibrary{automata, positioning, arrows.meta}
\usepackage[linesnumbered,ruled,vlined]{algorithm2e}



% ==============================
%         СОКРАЩЕНИЯ
% ==============================
\newcommand{\T}{\mathcal{T}}
\newcommand{\w}{\omega}
\newcommand{\pair}[2]{\langle #1, #2 \rangle}

% ==============================
%         ТИТУЛЬНЫЙ ЛИСТ
% ==============================
\begin{document}

\begin{titlepage}
\centering

{\LARGE \textbf{Теория формальных языков} \par}
\vspace{0.5cm}

{\large \textbf{Лабораторная работа №1} \par}
{\large Вариант 8 \par}
\vspace{1.5cm}

{\large \textbf{Работу выполнила:} \par}
{\large Студент группы ИУ9-51Б \par}
{\large Григорьева Софья \par}

\end{titlepage}


% ==============================
%         СОДЕРЖАНИЕ
% ==============================

\tableofcontents
\newpage

% ==============================
%           ЗАДАНИЕ
% ==============================

\section{Задание}

Дана SRS:
\begin{flushleft}
$apbc \rightarrow caqdbapbap$\\
$paqd \rightarrow daqdbapbap$\\
$ccpp \rightarrow adaqdqa$\\
$dpd \rightarrow pdp$\\
\end{flushleft}

По имеющейся SRS определить:
\begin{itemize}
    \item[---] завершимость;
    \item[---] конечность классов эквивалентности по НФ (для построения 
    эквивалентностей считаем, что правила могут применяться в обе стороны). 
    Если их конечное число, то построить минимальную систему переписывания, 
    соответствующую им;
    \item[---] локальную конфлюэнтность и пополняемость по Кнуту–Бендиксу.
\end{itemize}

По SRS $\mathcal{T}$ тем самым (исключая случай, когда она сразу локально 
конфлюэнтна или конечна и минимальна) строится другая SRS $\T'$, которая 
должна сохранять те же классы эквивалентности. Если исходная SRS завершима, 
то правила в $\T'$ должны удовлетворять условию убывания левой части 
относительно правой по выбранному фундированному порядку $\preceq$.

Провести автоматическое тестирование предполагаемой эквивалентности указанных SRS.\\\\
\textbf{Фазз-тестирование эквивалентности:} строится случайное слово $\w$ и 
случайная цепочка переписываний его в $\w'$ по $\T$. Проверить, можно ли 
получить $\w'$ из $\w$ (или наоборот) в рамках правил $\T'$.\\\\
\textbf{Метаморфное тестирование:} выбрать инварианты, которые должны сохраняться
(либо монотонно изменяться) при переписывании в рамках $\T$. Породить случайную 
цепочку переписываний над случайным словом в $\T'$ и проверить выполнение 
инвариантов. Как минимум два разных инварианта.
% -------------------------------------------------------------------------------------

\clearpage 
\section{Исследование SRS}



% ==============================
%        ЗАВЕРШИМОСТЬ
% ==============================
\subsection{Завершимость}

Исходная SRS $\mathcal{T}$ является завершаемой. \\

\noindent\textbf{Доказательство:} \\

Рассмотрим инвариант
\[
I(S) = \sum_{i=1}^{n} g(s_i) \, 2^{i-1}, \quad \text{где} \quad
g(x) = 
\begin{cases} 
10, & \text{если } x = c,\\ 
1, & \text{если } x = d,\\ 
0, & \text{иначе} 
\end{cases}
\]

Здесь $s_i$ — буква слова на $i$-й позиции (нумерация начинается с 1). \\

Проверим инвариант для правил переписывания.  
Пусть правило применяется к подстроке, первый символ которой находится на $m$-й позиции в строке ($1 \le m \le n$). \\

\begin{enumerate}[label=\Roman*.]
\item $apbc \rightarrow caqdbapbap$ \\\\
    $I(S') = I(S) - 62 \cdot 2^{\,m-1}$
\item $paqd \rightarrow daqdbapbap$ \\\\
    $I(S') = I(S) + 2^{\,m-1}$
\item $ccpp \rightarrow adaqdqa$ \\\\
    $I(S') = I(S) - 12 \cdot 2^{\,m-1}$
\item $dpd \rightarrow pdp$ \\
    $I(S') = I(S) - 3 \cdot 2^{\,m-1}$ \\
\end{enumerate}

Получаем убывающий инвариант для правил I, III и IV.



% ==============================
%     КЛАССЫ ЭКВИВАЛЕНТНОСТИ
% ==============================
\subsection{Классы эквивалентности}

Исходная SRS $\mathcal{T}$ имеет бесконечное количество классов эквивалентности. Например, строка вида $a^n$ образует отдельный класс эквивалентности для каждого $n \in \mathbb{N}$.



% ==============================
%    ЛОКАЛЬНАЯ КОНФЛЮЭНТНОСТЬ
% ==============================
\subsection{Локальная конфлюэнтность}

SRS $\T$ не является локально конфлюэнтной, так как слово $apbccpp$ за один шаг
может быть преобразовано в слова $caqdbapbapcpp$ (по правилу $apbc \rightarrow caqdbapbap$) 
и $apbadaqdqa$ (по правилу $ccpp \rightarrow adaqdqa$), но эти слова являются различными НФ.



% ==============================
%           ПОПОЛНЕНИЕ
% ==============================
\subsection{Пополняемость по Кнуту–Бендиксу}

Для применения алгоритма Кнута–Бендикса введём \emph{фундированный порядок} --- 
length-lexicographic order $\prec$ --- и, в соответствии с ним, переориентируем правила. 
(Такое преобразование сохраняет классы эквивалентности исходной SRS $\mathcal{T}$.) \\\\
После переориентации система переписываний примет вид:

\begin{flushleft} 
$caqdbapbap \rightarrow apbc$\\ 
$daqdbapbap \rightarrow paqd$\\ 
$adaqdqa \rightarrow ccpp$\\ 
$pdp \rightarrow dpd$ \\ 
\end{flushleft}

\noindent\textbf{Пополнение системы.} \\\\
\noindent\textbf{1-я итерация.} \\

\begin{tikzpicture}[node distance=2.5cm, auto]
  \node (A) {$caqdbapbapdp$};
  \node (B) [below left of=A] {$apbcdp$};
  \node (C) [below right of=A] {$caqdbapbadpd$};
  \node (R) [right of=C, xshift=4cm, align=center] {
    Добавляем правило: \\
    $caqdbapbadpd \rightarrow apbcdp$
  };
  \draw[->] (A) -- (B);
  \draw[->] (A) -- (C);
\end{tikzpicture} \\\\

\begin{tikzpicture}[node distance=2.5cm, auto]
  \node (A) {$daqdbapbapdp$};
  \node (B) [below left of=A] {$paqddp$};
  \node (C) [below right of=A] {$daqdbapbadpd$};
  \node (R) [right of=C, xshift=4cm, align=center] {
    Добавляем правило: \\
    $daqdbapbadpd \rightarrow paqddp$
  };
  \draw[->] (A) -- (B);
  \draw[->] (A) -- (C);
\end{tikzpicture}  \\\\

\begin{tikzpicture}[node distance=2.5cm, auto]
  \node (A) {$pdpdp$};
  \node (B) [below left of=A] {$dpddp$};
  \node (C) [below right of=A] {$pddpd$};
  \node (R) [right of=C, xshift=4cm, align=center] {
    Добавляем правило: \\
    $pddpd \rightarrow dpddp$
  };
  \draw[->] (A) -- (B);
  \draw[->] (A) -- (C);
\end{tikzpicture}  \\\\

Заметим, что на данном этапе ни одно правило нельзя редуцировать. \\\\

\textbf{2-я итерация.} \\\\

\begin{tikzpicture}[node distance=3cm, auto]
  \node (A) {$caqdbapbadpdaqdbapbadpd$};
  \node (B) [below left of=A] {$apbcdpaqdbapbadpd$};
  \node (C) [below right of=A] {$caqdbapbadppaqddp$};
  \node (R) [right of=C, xshift=1.5cm, align=center] {
    Добавляем правило: \\
    $caqdbapbadppaqddp \rightarrow$ \\
    $apbcdpaqdbapbadpd$
  };
  \draw[->] (A) -- (B);
  \draw[->] (A) -- (C);
\end{tikzpicture}  \\\\

Для практических целей остановимся на разумной длине правил.  
Добавим в систему правила из первой итерации, поскольку дальнейшее пополнение приводит к росту длины слов, а наша цель --- построить эквивалентную SRS, удобную для применения. Кроме того, можно доказать, что система бесконечно пополняется. \\\\

\noindent\textbf{Доказательство:} \\

Пусть в системе существует правило, левая часть которого $\sigma$ оканчивается на букву $p$, но не на $pp$, а правая часть $\sigma'$ не оканчивается на $p$. Обозначим через $\lambda$ префикс слова $\sigma$, исключая последнюю букву $p$. Тогда такое слово образует критическую пару с правилом $pdp$. \hfill(*) \\

\noindent\textbf{1-я итерация.}  \\

\begin{tikzpicture}[node distance=2.5cm, auto]
  \node (A) {$\sigma dp$};
  \node (B) [below left of=A] {$\sigma' dp$};
  \node (C) [below right of=A] {$\lambda dpd$};
  \draw[->] (A) -- (B);
  \draw[->] (A) -- (C);
\end{tikzpicture}  \\\\

Tак как правила системы либо уменьшают длину слова, либо сохраняют её, но уменьшают лексикографический порядок $\lambda \succ \sigma'$ и $\lambda dpd \succ \sigma' dp$, поэтому в систему добавляется правило $\lambda dpd \rightarrow \sigma' dp$. \\

Новое правило снова образует критическую пару с $pdp$. \hfill(*) \\

\noindent\textbf{2-я итерация.}  \\

\begin{tikzpicture}[node distance=2.5cm, auto]
  \node (A) {$\lambda dpdp$};
  \node (B) [below left of=A] {$\sigma' dpp$};
  \node (C) [below right of=A] {$\lambda ddpd$};
  \draw[->] (A) -- (B);
  \draw[->] (A) -- (C);
\end{tikzpicture}  \\\\

Аналогично, в систему добавляется правило $\lambda ddpd \rightarrow \sigma' dpp$ и оно снова образует критическую пару с $pdp$. \hfill(*) \\

Повторяя этот процесс, получаем бесконечную цепочку новых правил. 

(*) Образование новых критических пар и правил объясняется тем, что правила системы не могут редуцировать $dp^{+}$ на конце слова, а также последовательность букв $d$, если перед ней не стоит буква $p$, поэтому результат применения правил на схемах - различные нормальные формы.  

В исходной системе присутствует правило $caqdbapbap \rightarrow apbc$, левая часть которого оканчивается на $p$, но не на $pp$, а правая часть не оканчивается на $p$. Следовательно, система SRS $\tilde{\T}$ бесконечно пополняется.


% ==============================
%       ЭКВИВАЛЕНТНАЯ SRS
% ==============================
\subsection{Построение SRS $\tilde{\T}'$}

На основе исследования SRS $\tilde{\T}$ была построена эквивалентная ей SRS $\tilde{\T}'$:

\begin{flushleft} 
$caqdbapbap \rightarrow apbc$\\ 
$daqdbapbap \rightarrow paqd$\\ 
$adaqdqa \rightarrow ccpp$\\ 
$pdp \rightarrow dpd$ \\ 
$caqdbapbadpd \rightarrow apbcdp$ \\
$daqdbapbadpd \rightarrow paqddp$ \\
$pddpd \rightarrow dpddp$ \\
\end{flushleft}

% -------------------------------------------------------------------------------------
\newpage
\section{Тестирование}

% ==============================
%       ФАЗЗ-ТЕСТИРОВАНИЕ
% ==============================
\subsection{Фазз-тестирование эквивалентности}

Исходный код программы представлен в файле <<fuzzTest.erl>>. \\ 

В программе реализованы следующие основные функции:

\begin{itemize}[parsep=-1pt]
    \item \verb|originalSRS()| – возвращает список исходных правил;
    \item \verb|equalSRS()| – возвращает список правил, эквивалентных исходным;
    \item \verb|generateRandomWord()| – генерирует случайное слово длиной от 10 до 30 символов из алфавита \verb|"abcdpq"|;
    \item \verb|findOccurrences(Sub, Word)| – находит все позиции вхождений подстроки \verb|Sub| в строке \verb|Word|;
    \item \verb|randomlyTransform(Word, Rules)| – случайным образом применяет правила к слову;
    \item \verb|getNext(Word, Rules)| – возвращает все возможные варианты слов, которые могут быть получены из \verb|Word| применением одного правила;
    \item \verb|findTransformOptions(StartWord, Rules)| – находит все слова, достижимые из \verb|StartWord| при последовательном применении правил;
    \item \verb|areEqual(Small, Big)| – проверяет эквивалентность двух слов;
    \item \verb|test(TestsNumber)| – выполняет указанное количество фазз-тестов и выводит их результаты.
\end{itemize}

Программа последовательно генерирует случайные слова, применяет к ним исходные правила, и проверяет эквивалентность слов, через нахождение совпадений между результатами применения эквивалентной системы правил к этим словам, и выводит отчет о каждом тесте.

% ==============================
%    МЕТАМОРФНОЕ ТЕСТИРОВАНИЕ
% ==============================
\subsection{Метаморфное тестирование}

В качестве инвариантов для метаморфного тестирования были выбраны:

\begin{enumerate}[parsep=-1pt]
    \item сохранение чётности количества букв $c$
    \item сохранение чётности суммы количеств букв $b$ и $q$
\end{enumerate}

Исходный код программы представлен в файле <<metaTest.erl>>.\\  
В программе реализованы следующие основные функции:

\begin{itemize}[parsep=-1pt]
    \item \verb|originalSRS()| – возвращает список исходных правил;
    \item \verb|equalSRS()| – возвращает список правил, эквивалентных исходным;
    \item \verb|generateRandomWord()| – генерирует случайное слово длиной от 10 до 30 символов из алфавита \verb|"abcdpq"|;
    \item \verb|findOccurrences(Sub, Word)| – находит все позиции вхождений подстроки \verb|Sub| в строке \verb|Word|;
    \item \verb|randomlyTransform(Word, Rules)| – случайным образом применяет правила к слову;
    \item \verb|isValidInvarian1(Word, WordAfterOriginal, WordAfterEqual)| – проверяет сохранение первого инварианта (чётность количества букв $c$);
    \item \verb|isValidInvarian2(Word, WordAfterOriginal, WordAfterEqual)| – проверяет сохранение второго инварианта (чётность суммы количеств букв $b$ и $q$);
    \item \verb|test(TestsNumber)| – выполняет указанное количество метаморфных тестов и выводит результаты.
\end{itemize}

Программа генерирует случайные слова, применяет к ним исходную и эквивалентную системы правил, а затем проверяет, сохраняются ли заданные инварианты после преобразований и выводит отчет о каждом тесте.

\end{document}

